\documentclass[11pt]{article}

\usepackage[a4paper,margin=1in]{geometry}
\usepackage[T1]{fontenc}
\usepackage[utf8]{inputenc}
\usepackage{amsmath}
\usepackage{amssymb}
\usepackage{booktabs}
\usepackage{xcolor}
\usepackage{listings}
\usepackage{hyperref}

\hypersetup{
  colorlinks=true,
  linkcolor=blue!60!black,
  urlcolor=blue!60!black,
  citecolor=blue!60!black
}

\lstdefinestyle{code}{
  basicstyle=\ttfamily\small,
  breaklines=true,
  columns=fullflexible,
  frame=single,
  rulecolor=\color{black!20},
  backgroundcolor=\color{black!3},
  keepspaces=true
}
\lstset{style=code}

\title{Aero Quest Sim: Principles, Coordinate Frames, and Setup Tutorial}
\author{}
\date{}

\begin{document}
\maketitle

\section{Goal}

This project connects Meta Quest hand tracking to a MuJoCo simulation of an
SO101 arm with an Aero Hand. The current control path is a two-channel
controller:

\begin{lstlisting}
Quest wrist position        -> SO101 shoulder/elbow position IK
Quest palm direction        -> SO101 wrist_flex + wrist_roll
Quest wrist-local landmarks -> Aero Hand 7D finger action
\end{lstlisting}

The main entry point is:

\begin{lstlisting}
scripts/quest_so101_aero_nullspace_ik_teleop.py
\end{lstlisting}

Core files:

\begin{itemize}
  \item \texttt{aero\_quest/quest\_hand\_frame.py}: typed Quest frames, coordinate conversion, and palm-frame helpers.
  \item \texttt{aero\_quest/arm\_teleop.py}: task-space velocity control and damped least-squares IK.
  \item \texttt{aero\_quest/retargeting.py}: Quest 21-landmark to Aero Hand 7D action retargeting.
  \item \texttt{aero\_quest/so101\_aero\_control.py}: normalized action to MuJoCo actuator control mapping.
\end{itemize}

\section{Running the Demo}

Run from the project root:

\begin{lstlisting}
conda activate aero_sim
cd ~/Projects/aero_quest_sim
adb devices
adb reverse tcp:8000 tcp:8000
\end{lstlisting}

Quest Hand Tracking Streamer settings:

\begin{lstlisting}
Protocol: TCP
IP/Host: localhost or 127.0.0.1
Port: 8000
Hand: Right
\end{lstlisting}

Start full SO101 plus Aero Hand teleoperation:

\begin{lstlisting}
python scripts/quest_so101_aero_nullspace_ik_teleop.py
\end{lstlisting}

Useful MuJoCo viewer keys:

\begin{lstlisting}
R  Re-zero: map the current Quest wrist pose to the current SO101 end-effector pose
P  Pause or resume arm motion
\end{lstlisting}

\section{Input Data Format}

Quest TCP packets are converted into \texttt{QuestHandFrame} objects:

\begin{lstlisting}[language=Python]
hand_side: str
timestamp_ns: int | None
frame_id: int | None
wrist_pos_world: np.ndarray    # shape (3,)
wrist_quat_world: np.ndarray   # shape (4,), xyzw
landmarks_wrist: np.ndarray    # shape (21, 3)
\end{lstlisting}

The most important detail is that one packet contains data from two different
frames:

\begin{itemize}
  \item \texttt{wrist\_pos\_world} and \texttt{wrist\_quat\_world} are expressed in the Quest/Unity world frame $Q$.
  \item \texttt{landmarks\_wrist} are expressed in the Quest wrist/root local frame $Wrist$.
\end{itemize}

Do not treat \texttt{landmarks\_wrist} as world or robot coordinates. Convert
them to world coordinates only for visualization, debugging, or dataset
inspection:

\begin{lstlisting}[language=Python]
landmarks_Q = convert_landmarks_wrist_to_world(
    wrist_pos_world,
    wrist_quat_world,
    landmarks_wrist,
)
\end{lstlisting}

\section{Coordinate Frames}

\subsection{Quest World Frame $Q$}

The current Quest/Unity streamer uses:

\begin{lstlisting}
+X right
+Y up
+Z forward
\end{lstlisting}

\subsection{Robot Base Frame $B$}

The SO101/MuJoCo debug base frame used by this project is:

\begin{lstlisting}
+X forward
+Y left
+Z up
\end{lstlisting}

\subsection{Quest Wrist Frame $Wrist$}

\texttt{landmarks\_wrist} contains 21 points relative to the Quest hand
root/wrist. These points are best used for hand shape: finger curl, thumb
abduction, pinch distance, and similar features. They should not directly set
the robot end-effector position.

\section{\texorpdfstring{$R_{BQ}$}{R\_BQ}: Quest-to-Robot Axis Mapping}

The default matrix in the main script is:

\begin{lstlisting}[language=Python]
R_BQ = np.array([
    [ 0.0, 0.0, 1.0],
    [-1.0, 0.0, 0.0],
    [ 0.0, 1.0, 0.0],
])
\end{lstlisting}

It maps vectors from the Quest world frame $Q$ into the robot base frame $B$:

\begin{align}
\mathbf{v}_B = R_{BQ}\mathbf{v}_Q .
\end{align}

The default mapping is:

\begin{center}
\begin{tabular}{lll}
\toprule
Quest direction & Robot direction & Meaning \\
\midrule
$+Z_Q$ & $+X_B$ & hand forward moves the robot forward \\
$+X_Q$ & $-Y_B$ & hand right moves the robot right \\
$+Y_Q$ & $+Z_B$ & hand up moves the robot up \\
\bottomrule
\end{tabular}
\end{center}

If the arm moves in the wrong direction, tune this matrix first. The main
script can accept a custom matrix with \texttt{--R\_BQ} as 9 row-major floats.

\section{Teleop Zero and Relative Motion}

The Quest world origin is not treated as the robot base origin. At startup, or
when the user presses \texttt{R}, the controller records a teleop zero:

\begin{align}
\mathbf{p}^{0}_{wrist,Q} &= \text{current Quest wrist position}, \\
\mathbf{p}^{0}_{ee,B} &= \text{current SO101 end-effector position}, \\
R^{0}_{ee,B} &= \text{current SO101 end-effector orientation}.
\end{align}

After that, each frame uses only relative wrist displacement:

\begin{align}
\Delta\mathbf{p}_{Q}
  &= \mathbf{p}^{t}_{wrist,Q} - \mathbf{p}^{0}_{wrist,Q}, \\
\mathbf{p}^{target}_{B}
  &= \mathbf{p}^{0}_{ee,B}
     + s R_{BQ}\Delta\mathbf{p}_{Q}.
\end{align}

Here $s$ is the position scale. The main script default is:

\begin{lstlisting}
--scale 0.6
\end{lstlisting}

This means 1 meter of Quest hand motion produces about 0.6 meters of robot
target motion. In practice, smaller values are usually easier to control.

\section{Workspace Limits}

The computed target position is clipped to an axis-aligned workspace box:

\begin{lstlisting}
--workspace-min 0.05 -0.35 0.03
--workspace-max 0.55  0.35 0.60
\end{lstlisting}

In robot-base coordinates:

\begin{align}
x_B &\in [0.05, 0.55], \\
y_B &\in [-0.35, 0.35], \\
z_B &\in [0.03, 0.60].
\end{align}

Interpretation:

\begin{itemize}
  \item $x_B$: forward/backward, larger means farther forward.
  \item $y_B$: left/right, larger means farther left.
  \item $z_B$: height, larger means higher.
\end{itemize}

If the end effector stops moving near a boundary, inspect \texttt{target\_B} in
the debug output and check whether it is being clipped by the workspace.

\section{Palm Direction and Wrist Control}

Palm direction comes from \texttt{landmarks\_wrist}. The relevant landmark
indices are:

\begin{lstlisting}
0  wrist
5  index MCP
9  middle MCP
17 pinky MCP
\end{lstlisting}

First, the code builds a palm frame in the $Wrist$ frame:

\begin{align}
\mathbf{x}_{palm}
  &= \operatorname{normalize}(\mathbf{P}_5-\mathbf{P}_{17}), \\
\mathbf{y}_{hint}
  &= \operatorname{normalize}(\mathbf{P}_9-\mathbf{P}_0), \\
\mathbf{z}_{palm}
  &= \operatorname{normalize}(\mathbf{x}_{palm}\times\mathbf{y}_{hint}), \\
\mathbf{y}_{palm}
  &= \operatorname{normalize}(\mathbf{z}_{palm}\times\mathbf{x}_{palm}).
\end{align}

Thus:

\begin{align}
R_{palm,Wrist}
  = [\mathbf{x}_{palm},\mathbf{y}_{palm},\mathbf{z}_{palm}] .
\end{align}

The Quest wrist quaternion then maps the palm frame into Quest world:

\begin{align}
R_{palm,Q} = R_{wrist,Q}R_{palm,Wrist}.
\end{align}

If orientation control is enabled, the controller computes a relative rotation:

\begin{align}
R_{\Delta,Q} &= R^{t}_{palm,Q}(R^{0}_{palm,Q})^T, \\
R_{\Delta,B} &= R_{BQ}R_{\Delta,Q}R_{BQ}^T, \\
R^{target}_{B} &= R_{\Delta,B}R^{0}_{ee,B}.
\end{align}

In the current full teleop script, position IK mainly controls
\texttt{shoulder\_pan}, \texttt{shoulder\_lift}, and \texttt{elbow\_flex}.
Palm direction is also projected onto \texttt{wrist\_flex} and
\texttt{wrist\_roll}:

\begin{lstlisting}
--wrist-flex-axis +y
--wrist-roll-axis -x
--wrist-flex-gain 1.0
--wrist-roll-gain 1.0
--wrist-orientation-alpha 0.25
--max-wrist-orientation-step 0.03
\end{lstlisting}

\texttt{wrist-orientation-alpha} controls how quickly the wrist follows the
target. \texttt{max-wrist-orientation-step} limits the target change per
simulation step and helps avoid sudden jumps.

\section{Task-Space Velocity Control}

The arm does not teleport to \(\mathbf{p}^{target}_{B}\). It first computes a
position error:

\begin{align}
\mathbf{e}_{p} = \mathbf{p}^{target}_{B} - \mathbf{p}^{current}_{B}.
\end{align}

Then it computes a bounded desired linear velocity:

\begin{align}
\mathbf{v}_{cmd}
  = \operatorname{clipNorm}(k_p\mathbf{e}_p, v_{max}).
\end{align}

Default values:

\begin{lstlisting}
--kp-pos 5.0
--max-linear-speed 0.18
\end{lstlisting}

A larger \texttt{kp-pos} makes the end effector follow the target more
aggressively. \texttt{max-linear-speed} is the linear speed cap.

\section{Damped Least-Squares IK}

IK uses MuJoCo to compute the Jacobian of the end-effector site, then selects
only the arm joints. The default arm joints are:

\begin{lstlisting}
shoulder_pan,shoulder_lift,elbow_flex
\end{lstlisting}

Let the task-space velocity be \(\dot{\mathbf{x}}\), the Jacobian be \(J\), and
the damping coefficient be \(\lambda\). The damped least-squares solution is:

\begin{align}
\dot{\mathbf{q}}
  = J^T(JJ^T+\lambda^2I)^{-1}\dot{\mathbf{x}}.
\end{align}

The solver integrates this velocity into a joint-position target:

\begin{align}
\mathbf{q}^{target}
  = \mathbf{q}^{current}+\dot{\mathbf{q}}\Delta t.
\end{align}

Relevant parameters:

\begin{lstlisting}
--ik-damping 0.05
--max-joint-speed 1.2
--joint-target-smoothing-alpha 0.0
\end{lstlisting}

Higher \texttt{ik-damping} is usually more stable but less responsive. Lower
values are faster but can shake near singular configurations.
\texttt{max-joint-speed} limits each joint velocity.

\section{Aero Hand 7D Finger Action}

The hand channel uses only \texttt{landmarks\_wrist}. It first converts the 21
points into a palm-local frame, removing global position, global orientation,
and approximate hand-size differences:

\begin{align}
\mathbf{P}_{local}
  = \frac{(\mathbf{P}-\mathbf{P}_0)R_{palm}}{\|\mathbf{P}_9-\mathbf{P}_0\|}.
\end{align}

The output is a 7D Aero Hand action:

\begin{center}
\begin{tabular}{cll}
\toprule
Index & Name & Meaning \\
\midrule
0 & thumb\_abduction & thumb abduction/adduction \\
1 & thumb\_flexion\_1 & first thumb flexion channel \\
2 & thumb\_flexion\_2 & second thumb flexion channel \\
3 & index\_curl & index finger curl \\
4 & middle\_curl & middle finger curl \\
5 & ring\_curl & ring finger curl \\
6 & little\_curl & little finger curl \\
\bottomrule
\end{tabular}
\end{center}

Each value is clamped to $[0,1]$. For finger curl:

\begin{lstlisting}
0 = open
1 = closed
\end{lstlisting}

For three points around a joint, the bend angle is:

\begin{align}
\theta
  &= \arccos
     \frac{(\mathbf{A}-\mathbf{B})\cdot(\mathbf{C}-\mathbf{B})}
          {\|\mathbf{A}-\mathbf{B}\|\|\mathbf{C}-\mathbf{B}\|}, \\
bend &= \pi-\theta.
\end{align}

The bend is normalized with:

\begin{align}
curl
  = \operatorname{clip}_{[0,1]}
    \frac{bend-open\_angle}{closed\_angle-open\_angle}.
\end{align}

Current defaults:

\begin{lstlisting}
open_angle   = 0.08 rad
closed_angle = 1.35 rad
\end{lstlisting}

For index, middle, ring, and little fingers, curl combines proximal and distal
bend:

\begin{align}
curl
  = normalize(0.65\,bend_{proximal}+0.35\,bend_{distal}).
\end{align}

\section{Mapping Actions to MuJoCo Actuators}

The SO101 plus Aero Hand actuator mapping lives in:

\begin{lstlisting}
aero_quest/so101_aero_control.py
\end{lstlisting}

SO101 arm actuator names:

\begin{lstlisting}
shoulder_pan
shoulder_lift
elbow_flex
wrist_flex
wrist_roll
\end{lstlisting}

Aero Hand semantic action mapping:

\begin{lstlisting}
thumb_abduction -> right_thumb_A_cmc_abd
thumb_flexion_1 -> right_th1_A_tendon
thumb_flexion_2 -> right_th2_A_tendon
index_curl      -> right_index_A_tendon
middle_curl     -> right_middle_A_tendon
ring_curl       -> right_ring_A_tendon
little_curl     -> right_pinky_A_tendon
\end{lstlisting}

The hand action is in $[0,1]^7$. When it is written into MuJoCo control values,
each channel is linearly mapped into the actuator \texttt{ctrlrange}. In the
current Aero Hand mapping, each channel is marked \texttt{inverted=True}, so the
code first applies:

\begin{align}
a' = 1-a
\end{align}

and then maps the result to the actuator range.

\section{Debugging Checklist}

If Quest does not connect, check:

\begin{lstlisting}
adb devices
adb reverse --list
\end{lstlisting}

If the arm moves in the wrong direction, inspect these values in order:

\begin{enumerate}
  \item \texttt{delta\_p\_Q}: whether Quest wrist motion is detected in the expected direction.
  \item \texttt{R\_BQ}: whether Quest axes map to the correct robot axes.
  \item \texttt{target\_B}: whether the target is being clipped by the workspace.
  \item \texttt{ee\_B} and \texttt{err}: whether IK can track the target.
\end{enumerate}

If finger motion looks unnatural, check:

\begin{enumerate}
  \item Whether \texttt{landmarks\_wrist} has shape \((21,3)\).
  \item \texttt{open\_angle} and \texttt{closed\_angle} in \texttt{aero\_quest/retargeting.py}.
  \item \texttt{--hand-smoothing-alpha}: too small is sluggish; too large can shake.
\end{enumerate}

\section{Recommended Learning Order}

\begin{enumerate}
  \item Run \texttt{scripts/debug\_quest\_dual\_channel.py} to verify wrist pose and landmarks.
  \item Run \texttt{scripts/quest\_arm\_channel\_target\_ball.py} to inspect translation axes only.
  \item Run \texttt{scripts/quest\_arm\_channel\_so101\_ik.py} to tune the SO101 arm channel only.
  \item Run \texttt{scripts/quest\_tcp\_aero\_teleop.py --alpha 0.25} to tune Aero Hand finger retargeting only.
  \item Run \texttt{scripts/quest\_arm\_channel\_so101\_aero\_full\_teleop.py} to combine both channels.
\end{enumerate}

\section{Core Rules}

The most important design rule is to keep coordinate frames and control
channels explicit:

\begin{itemize}
  \item The Quest world origin is not the robot base origin.
  \item Quest wrist world pose is used by the arm channel.
  \item Wrist-local landmarks are used by the hand channel.
  \item Robot target position is based on relative motion from teleop zero.
  \item Axis differences are handled by $R_{BQ}$.
  \item Workspace bounds, velocity caps, and IK damping are the main stability safeguards.
\end{itemize}

\end{document}
